%% C-rationale.tex — Annex C: Language design rationale (informative)
\chapter{Language design rationale}
\label{annex:C-rationale}
\index{rationale}

\pnum
This annex is informative.  It explains the reasoning behind key design
decisions in the Lain language.  Understanding the rationale helps readers
interpret edge cases in the normative text and guides implementors when
the specification is ambiguous.

% =========================================================================
\section{Declaring effects: the \texttt{effects} row}
\label{sec:rationale.func-proc}
\index{effects row!rationale}
% =========================================================================

\pnum
\textbf{Decision.}  State what a function may do at the declaration level, in an
\lain{effects} row on the one introducer \lain{func}, where silence means the
empty row --- pure and total.

\pnum
\textbf{Rationale.}  Every reason below was a reason to declare effects at all,
and each still holds; only the spelling changed.
\begin{itemize}
  \item \textit{Termination guarantees}: a function whose row does not name
        \lain{diverge} is guaranteed to terminate (no unbounded \lain{while}, no
        unchecked recursion).  This lets the implementation reason statically
        about completion, and lets such a body be used inside \lain{comptime}
        expressions.
  \item \textit{Effect propagation}: the constraints propagate --- a function may
        not call one whose row names an effect its own row does
        not~(\diagcode{E011}) --- so pure code stays pure by construction.
  \item \textit{Optimisation}: the implementation may freely reorder, memoize or
        eliminate calls to a function with the empty row, because it has no
        observable side effects.
  \item \textit{Readability}: the row tells the reader what a callable may do
        without inspecting its body, and because it is a complete upper bound
        rather than a hint, the \emph{absence} of a row is itself a statement.
\end{itemize}

\pnum
\textbf{The original spelling, and why it was replaced.}  The decision was first
realised as two introducers: \lain{func} for the pure and total case and
\lain{proc} for anything else.  That form was removed on 2026-09-25, in the
effect-row work (plan Part~7B).  Three reasons, each about what a keyword
\emph{cannot} say:
\begin{itemize}
  \item The bound had only two points, \lain{func} and \lain{proc} --- a single
        bit.  Four effects therefore had to be spelled a second way regardless,
        and were: by keyword, by attribute (\lain{@diverges}) and by clause, all
        at once.  One concern with three spellings is what the language's own law
        of a single mechanism refuses.
  \item A row makes assignment between function-pointer types \emph{row
        containment}, which subsumes the one relation the keywords could express
        (a \lain{func} being usable where a \lain{proc} was expected) and gives
        every other one for free.
  \item A row lets an arrow say \emph{``may print, but terminates''}.  A bit
        cannot: it must either grant divergence along with \lain{io} or deny
        both.
\end{itemize}

\pnum
\textbf{Alternatives considered.}
A single keyword with a purity annotation (\lain{func(pure)}) was rejected as less
visible than a row that must name every effect.  A Haskell-style IO monad was
rejected as too complex for a systems language.

% =========================================================================
\section{Explicit \texttt{mov} at call sites}
\label{sec:rationale.mov}
\index{mov!rationale}
% =========================================================================

\pnum
\textbf{Decision.}  Require explicit \lain{mov} at call sites when
transferring ownership.

\pnum
\textbf{Rationale.}
\begin{itemize}
  \item \textit{Readability}: \lain{close\_file(mov f)} makes ownership
        transfer visible at the call site.
  \item \textit{Intentionality}: accidental ownership transfer is a
        common source of bugs.  Requiring \lain{mov} forces the
        programmer to acknowledge the transfer.
  \item \textit{Symmetry}: the annotation appears in both declarations
        (\lain{mov x T}) and at call sites (\lain{mov x}), creating a
        uniform language.
\end{itemize}

\pnum
\textbf{Alternatives considered.}
Rust-style implicit moves were rejected because they are invisible at
call sites.  A library function \texttt{std::move} (C++) was rejected
because ownership transfer should be a language primitive, not a library
function.

% =========================================================================
\section{Implicit NLL lifetimes (no annotations)}
\label{sec:rationale.nll}
\index{lifetime annotations!rationale}
% =========================================================================

\pnum
\textbf{Decision.}  Borrows have implicit lifetimes determined by
non-lexical lifetime analysis.  No explicit lifetime annotations are
required.

\pnum
\textbf{Rationale.}
\begin{itemize}
  \item \textit{Simplicity}: Rust's lifetime annotations are widely
        cited as its steepest learning curve.  Lain targets programmers
        who want safety without that complexity.
  \item \textit{Sufficient expressiveness}: NLL handles the vast
        majority of real-world borrow patterns.
  \item \textit{Predictability}: borrows expire at last use — a single,
        intuitive rule.
\end{itemize}

\pnum
\textbf{Trade-offs.}  Some programs that Rust can verify with explicit
lifetime annotations may be rejected by Lain's more conservative
analysis.  Higher-order borrowing patterns are constrained.

% =========================================================================
\section{C99 translation output}
\label{sec:rationale.c99}
\index{C99 backend!rationale}
% =========================================================================

\pnum
\textbf{Decision.}  Compile to C99 source code rather than LLVM IR,
machine code, or bytecode.

\pnum
\textbf{Rationale.}
\begin{itemize}
  \item \textit{Portability}: C99 compilers exist for every platform,
        from server-class x86 to microcontrollers.  Lain inherits this
        portability for free.
  \item \textit{Toolchain reuse}: users get GCC/Clang optimisation passes,
        debuggers, profilers, and sanitisers without custom tooling.
  \item \textit{Interoperability}: C interop is zero-cost because the
        output \emph{is} C.
  \item \textit{Simplicity}: a C emitter is orders of magnitude simpler
        than an LLVM backend or a native code generator.
  \item \textit{Auditability}: the generated C code is human-readable.
\end{itemize}

\pnum
\textbf{Trade-offs.}  Compilation is two-phase (Lain~→~C~→~binary).
Some optimisations are impossible without direct IR control.  The C type
system constrains what Lain can express (e.g., no guaranteed tail calls).

% =========================================================================
\section{No exception mechanism}
\label{sec:rationale.no-exceptions}
\index{exceptions!absence of, rationale}
% =========================================================================

\pnum
\textbf{Decision.}  Lain has no exception mechanism.  All error handling
uses return values.

\pnum
\textbf{Rationale.}
\begin{itemize}
  \item \textit{Zero overhead}: no unwinding tables, catch tables, or
        RTTI.
  \item \textit{Determinism}: error handling follows the same control
        flow as normal code; no invisible jumps.
  \item \textit{Embedded compatibility}: stack unwinding is unavailable
        or unreliable on many embedded platforms.
  \item \textit{Explicit error paths}: every error is visible in the
        function signature, as a marker of its union return type
        (\lain{T | NotFound}, \S\,\ref{sec:types.union}).
  \item \textit{Compatibility with \lain{defer}}: the \lain{defer}
        mechanism provides deterministic cleanup without exceptions.
\end{itemize}

% =========================================================================
\section{Explicit \texttt{unsafe} blocks}
\label{sec:rationale.unsafe}
\index{unsafe!rationale}
% =========================================================================

\pnum
\textbf{Decision.}  Pointer dereference and direct ADT field access
require an explicit \lain{unsafe} block.

\pnum
\textbf{Rationale.}
\begin{itemize}
  \item \textit{Auditability}: safety-critical codebases can grep for
        \lain{unsafe \{} to enumerate all potentially dangerous
        operations.
  \item \textit{Minimal scope}: Lain's \lain{unsafe} is narrow — it
        enables only pointer dereference, address-of, casts, direct ADT
        access, and bounds suppression.  Unlike C (where everything is
        unsafe by default) or some languages (where \lain{unsafe} enables
        a much larger capability set), Lain's \lain{unsafe} is a
        surgical escape hatch.
  \item \textit{Social contract}: the keyword communicates to code
        reviewers that extra scrutiny is warranted.
\end{itemize}

% =========================================================================
\section{Value Range Analysis instead of SMT}
\label{sec:rationale.vra}
\index{VRA!rationale}
% =========================================================================

\pnum
\textbf{Decision.}  Use a value range analysis --- intervals, and the relations between
pairs of variables an octagon keeps --- for static verification instead of an SMT solver.

\pnum
\textbf{Rationale.}
\begin{itemize}
  \item \textit{Decidability}: VRA always terminates in polynomial time.
        SMT solving is NP-hard in the general case and may time out.
  \item \textit{Predictability}: programmers can reason about what the
        compiler will accept by thinking about integer ranges.  SMT
        solver reasoning is opaque.
  \item \textit{Self-containment}: the analysis is part of the compiler, with no
        external dependency, and its every verdict can be replayed from a certificate.
\end{itemize}

\pnum
\textbf{Trade-offs.}  VRA is incomplete — it may reject valid programs
that an SMT solver could prove correct.  Non-linear constraints produce
conservatively wide ranges.

% =========================================================================
\section{Type parameters instead of \texttt{<T>} syntax}
\label{sec:rationale.comptime}
\index{generics!rationale}
% =========================================================================

\pnum
\textbf{Decision.}  A type parameter is an ordinary parameter whose type is \lain{type},
\lain{func max(T type, a T, b T) T}, and every instance is monomorphized
(\S\,\ref{sec:14-generics}).

\pnum
\textbf{Rationale.}
\begin{itemize}
  \item \textit{Unification}: types and compile-time values are passed the same way;
        \lain{type Buf(N usize)} takes a length as \lain{type Box(T type)} takes a type.
  \item \textit{Simplicity}: no angle brackets, no trait bounds, no where clauses.  The
        body is checked at each instance with the types it is given.
  \item \textit{Zero cost}: an instance is an ordinary function or type, with no
        dictionary passing and no boxing.
  \item \textit{Precedent}: inspired by Zig, whose types are compile-time values.
\end{itemize}

% =========================================================================
\section{Universal Function Call Syntax (UFCS)}
\label{sec:rationale.ufcs}
\index{UFCS!rationale}
% =========================================================================

\pnum
\textbf{Decision.}  Support UFCS: \lain{x.f(args)} desugars to
\lain{f(x, args)}.

\pnum
\textbf{Rationale.}
\begin{itemize}
  \item \textit{No classes needed}: UFCS provides method-like syntax
        without class hierarchies, virtual dispatch, or inheritance.
  \item \textit{Chainability}: enables fluent APIs.
  \item \textit{Discoverability}: IDE auto-completion works naturally —
        typing a dot after a value shows available operations.
  \item \textit{Simplicity}: syntactic sugar only; no new concepts, no
        method tables, no implicit \lain{this} pointer.
\end{itemize}

% =========================================================================
\section{Bounded \texttt{while} in \texttt{func}}
\label{sec:rationale.bounded-while}
\index{bounded while!rationale}
% =========================================================================

\pnum
\textbf{Decision.}  Allow \lain{while} loops in a function with the empty row when
their termination can be shown: by a measure the compiler infers from the loop, or by one
the programmer writes with \lain{decreasing}.

\pnum
\textbf{Rationale.}
A \textit{measure} --- a non-negative integer expression that strictly decreases on every
iteration --- is a well-established technique from program verification
(\textit{ranking functions}).  For the common loop, \lain{while i < n} with \lain{i}
rising, the compiler finds it itself (\S\,\ref{sec:stmt.while.infer}); where it cannot, the
programmer states one, and the compiler checks it rather than believing it.  Neither needs
an external solver.

\pnum
This allows algorithms such as binary search and GCD computation to be
expressed with the empty effect row, rather than having to admit an effect they do not have.
